\section{Methodology}
We selected two advanced reactors for this work: the X-energy Xe-100 and 
the \gls{USNC} \gls{MMR}. These two reactors provide comparison across 
some notable reactor design spaces, such as power output and fuel form. 
Additionally, there is enough information about each design publicly 
available to model each design with a reasonable degree of accuracy. 
we created models of each core and performed neutronics analysis using 
Serpent \cite{leppanen_serpent_2014}. 

We used the Xe-100-like model from \cite{richter_phlox_2021}, which can 
be found at \cite{richter_zoerichterphlox_2022}. 

We created a model for the \gls{MMR} primarily based on information in 
\cite{hawari_development_2018}, and supplemented or modified based on 
information published by \gls{USNC} \cite{noauthor_usnc_2021}. Figure 
\ref{fig:mmr_core} shows the radial and axial geometry of the \gls{MMR} 
model. 

\begin{figure}
        \begin{subfigure}{0.45\textwidth}
                \includegraphics[scale=0.07]{bachmann-mmr_geom1.png}
                \caption{Radial view of the USNC MMR model.}
                \label{fig:mmr_radial}
        \end{subfigure}
        \hfill 
        \begin{subfigure}{0.45\textwidth}
                \includegraphics[scale=0.1]{bachmann-mmr_geom3.png}
                \caption{Radial view of the USNC MMR model.}
                \label{fig:mmr_axial}
        \end{subfigure}
        \caption{Geometry of USNC MMR core model in Serpent.}
        \label{fig:mmr_core}
\end{figure}
